\section{Uniform group-ring coefficient quotients for the ramified Artin factor}
\label{sec:ray-coefficient}

Throughout this section \(p=269\), \(K=\mathbb Q(\sqrt{-3})\), and
\(\chi:G_K\to\mu _3\subset\mathcal O^\times\) is the fixed nontrivial cubic
character used in AC-K33. Fix \(b\in\{1,2\}\), put
\[
 \rho_0=\operatorname{Ind}_{G_K}^{G_{\mathbb Q}}\chi^b,\qquad
 D_0=D_{\rm dR}(\rho_{0,p}(1)),
\]
and fix an \(\mathcal O\)-lattice \(D_0^\circ\) in \(D_0\).
As in AC-K33, \(D_0\) is a two-dimensional multiplicity label with
trivial global Galois action. Let \(\Delta\) be a finite abelian
\(p\)-group quotient of \(G_K\) unramified at the unique prime over \(p\).
Write
\[
 A=\mathcal O[\Delta],\quad\vartheta_\Delta(g)=[\bar g],\quad
 \rho_\Delta=\operatorname{Ind}_{G_K}^{G_{\mathbb Q}}
                    (\chi^b\vartheta_\Delta).
 \tag{RCQ.1}
\]
This group ring is reduced and finite free over \(\mathcal O\); its
special fiber is generally nonreduced and its generic fiber generally
disconnected. Every dual below specifies its coefficient ring.

Published inputs are Colmez--Wang, \emph{Factorisation du syst\`eme de
Beilinson--Kato}, Theorem 13.11 and its complete printed proof,
\S13.4.1, pp.~129--130 in
\url{https://webusers.imj-prg.fr/~pierre.colmez/KE.pdf};
their defining intersection (3.30) and Theorem 3.34(i);
Helm--Moss, \emph{Converse theorems and the local Langlands
correspondence in families}, Invent.\ Math.\ 214 (2018), Lemma 3.1
and Corollary 7.7; and Helm's Theorem 7.9 in
\url{https://www.ma.imperial.ac.uk/~dhelm/bernstein3.pdf},
whose conjectural hypothesis the Helm--Moss corollary discharges.
The proof below does not reprove their deep upstream inputs.

\begin{lemma}[AC-K33-ray-period: an integral unramified period]
\label{AC-K33-ray-period}
For \(z\in\Delta\), there is an unramified extension \(L/\mathbb Q_p\)
and a unit
\[
 t_z\in\mathcal O_L\otimes_{\mathbb Z_p}A,\qquad
 \varphi(t_z)=z\,t_z.
 \tag{RCQ.2}
\]
Both the period and its inverse are integral, with operator norm at
most one in the coefficient lattices, uniformly in the order of \(z\).
\end{lemma}
\begin{proof}
Let \(n\) be the order of \(z\), take \(L/\mathbb Q_p\) unramified of
degree \(n\), and choose an integral normal-basis generator \(\alpha\).
Such a generator is obtained by lifting a normal-basis generator of
\(\mathbb F_{p^n}/\mathbb F_p\), since the determinant of its conjugates
is then a \(p\)-adic unit. Set
\(t_z=\sum_{i=0}^{n-1}\varphi^i(\alpha)z^{-i}\).
Cyclic reindexing proves (RCQ.2). The trace of the residue of \(\alpha\)
is nonzero: otherwise its conjugates have the nontrivial linear
relation with every coefficient one. Hence
\(\operatorname{Tr}_{L/\mathbb Q_p}\alpha\in\mathbb Z_p^\times\).
Modulo the maximal ideal of \(A\), every element of \(\Delta\) is one
and \(t_z\) is this unit trace. It is therefore a unit in the finite
\(A\)-algebra \(\mathcal O_L\otimes A\), even if that algebra splits.
Its inverse is integral. Using \(t_z^{-1}\) gives the opposite
Frobenius orientation with the same bounds.
\end{proof}

\begin{lemma}[AC-K33-ray-Schwartz: an integral compact vector on every component]
\label{AC-K33-ray-Schwartz}
Let \(\ell\ne p\), and let
\(r:G_{\mathbb Q_\ell}\to{\rm GL}_2(A)\) have finite image.
In the rational product of the classical Colmez--Wang Whittaker
models at the minimal primes of \(A\), the vector with Kirillov
restriction \(1_{\mathbb Z_\ell^\times}\) has integral values under
every \({\rm GL}_2(\mathbb Q_\ell)\)-translate, in the descended
coefficient convention of Colmez--Wang \S3.4.1. Its value at the
identity is one. No connectedness, constant-conductor assumption
or division by \(|\Delta|\) is needed.
\end{lemma}
\begin{proof}
Extend coefficients faithfully flatly to the completed maximal
unramified coefficient ring. The residue field is then algebraically
closed and the requisite additive character is available. The
base-changed \(A\) remains reduced, complete, Noetherian and flat
over its Witt ring. Helm--Moss Corollary 7.7, via Helm Theorem 7.9,
gives a co-Whittaker family. Normalize its top derivative to \(A\).
The composite of Helm--Moss (3.1), Lemma 3.1, is literally the natural
inclusion
\[
 {\rm cInd}_{U}^{P}\psi_A\longrightarrow V
 \longrightarrow{\rm Ind}_{U}^{G}\psi_A
 \longrightarrow{\rm Ind}_{U}^{P}\psi_A.
 \tag{RCQ.3}
\]
Apply it to the compact mirabolic function with restriction
\(1_{\mathbb Z_\ell^\times}(x)\) on \(\operatorname{diag}(x,1)\).
Its middle Whittaker function is integral on all \(G\).

The normalization matters. Emerton--Helm \S4.2, Definition 4.2.1,
uses \({\rm LL}(r)|\det|_\ell^{-1/2}\) for an irreducible generic
two-dimensional fiber. The family in Helm--Moss Conjecture 7.1(3)
has the dual of that construction for \(r^\vee\), hence has generic
fiber \({\rm LL}(r)|\det|_\ell^{+1/2}\).
Twisting it by \(|\det|_\ell^{-1}\) gives
\({\rm LL}(r)|\det|_\ell^{-1/2}\). Its fibers now have the Colmez--Wang
convention, with central character \(|\cdot|_\ell^{-1}\det r\);
in the split case this is the unnormalized
\({\rm Ind}_B^G(\chi_1\otimes|\cdot|_\ell^{-1}\chi_2)\)
(Colmez--Wang \S1.2.4 and \S3.4.4).
This unit-valued integral twist leaves our compact restriction
unchanged.

The finite Artin fibers cannot be a character times
\(1\oplus\epsilon_p\): a ratio of finite-order Frobenius eigenvalues
cannot equal \(\ell\) or \(\ell^{-1}\).
Their classical representations are irreducible generic, so the
Whittaker restriction to the mirabolic is injective. The function
just constructed is therefore the unique rational function with
the stated restriction on each generic component.

Descend without averaging using the twisted semilinear descent of
Emerton--Helm Proposition 3.1.4. Its proof uses faithful flatness and
\(\psi^\sigma(x)=\psi(e_\sigma x)\),
\(e_\sigma\in\mathbb Z_\ell^\times\).
Our restriction is fixed by \(x\mapsto e_\sigma x\).
Uniqueness on every generic component gives descent of the
Whittaker function. This is precisely the convention
\(\sigma_a\phi(x)=\phi(ax)\) in Colmez--Wang \S3.4.1.
Faithful-flat descent preserves integrality; (RCQ.3) gives value
one at \(x=1\).
\end{proof}

\begin{theorem}[AC-K33-ray-coefficient: a uniform actual coefficient quotient]
\label{AC-K33-ray-coefficient}
Fix \(j\ge0\). There are \(a\ge1\) and \(C\ge0\), depending only on
\(j,\rho_0,\mathcal O,D_0^\circ\), such that for every \(\Delta\)
in (RCQ.1) there is a principal tame level \(M_\Delta\), prime to \(p\),
and \(N_\Delta=M_\Delta p^a\), for which the actual inverse tower of
full \(\mathbb Q\)-schemes admits a continuous equivariant surjection
\[
 \begin{split}
 &\left(\varprojlim_r H^1_{\rm et}
  \bigl(Y(M_\Delta,N_\Delta p^r)_{\overline{\mathbb Q}},
                  {\rm TSym}^{j}\mathcal H(1)\bigr)\right)[1/p]\\
 &\hspace{24mm}\longtwoheadrightarrow
       \rho_\Delta^{\vee_A}[1/p]\otimes_{\mathcal O}D_0^\vee .
 \end{split}
 \tag{RCQ.4}
\]
The map is \(T\)-equivariant for the actual Hecke quotient \(T\to A\)
constructed below; equivalently it factors through the completed
\(A\)-quotient of its source.
The target inverse diamond is \(u^{-j-1}\), and \(U_p'=0\).
The integral source maps into
\[
 p^{-C}\bigl(\rho_\Delta^{\vee_A}\otimes_{\mathcal O}
                            (D_0^\circ)^\vee\bigr),
 \tag{RCQ.5}
\]
with \(C\) independent of the order and tame conductor of \(\Delta\).
The tame level may vary. No compatibility with geometric \(K\)-ray
norm maps, and no nonvanishing of a particular Rankin class, is
asserted here.
\end{theorem}
\begin{proof}
\textit{1. The actual group-ring Hecke quotient.}
Every characteristic-zero specialization of
\(\chi^b\vartheta_\Delta\) is finite order and differs from its
conjugate: equality would identify a nontrivial order-three
character with a character of \(p\)-power order.
Its theta series is cuspidal of weight one.
Its global and \(p\)-local residual representation equals the same
absolutely irreducible residual representation as \(\rho_0\).
Thus Pan's weight-one completed-cohomology input used in AC-K33
applies to every specialization. There are only finitely many
specializations, so a common principal tame level supports them.

Let \(T\) be the actual complete reduced local Hecke ring, with
free rank-two \(\rho_T\), of Colmez--Wang \S13.1.1--2.
Simultaneous evaluation gives a continuous map
\[
                         T\longrightarrow A.
 \tag{RCQ.6}
\]
Indeed the dense good-Hecke algebra has values in \(A\), equal to
the traces and determinants of \(\rho_\Delta\).
The image of \(A\hookrightarrow\prod_\psi\mathcal O_\psi\) is closed
because \(A\) is finite over \(\mathcal O\). Evaluation into this
finite product is continuous. Passing to its closure proves (RCQ.6),
without evaluating arbitrary formal power series at unrelated points.

The map is onto. Choose \(h\in G_K\) with
\(\chi^b(h)=\omega\) and
\(\vartheta_\Delta(h)=\vartheta_\Delta^c(h)=1\).
The order-three quotient is disjoint from the relevant finite
\(p\)-group quotient, so such an \(h\) exists; it need not have
order three in the absolute Galois group.
Writing \(t(g)=\operatorname{tr}\rho_\Delta(g)\), for \(g\in G_K\) we have
\[
 \chi^b(g)\vartheta_\Delta(g)=
       \frac{t(hg)-\omega^{-1}t(g)}{\omega-\omega^{-1}}.
 \tag{RCQ.7}
\]
The denominator is a unit. Continuity and Chebotarev put all these
Galois traces in the Hecke image, which consequently contains every
group-ring basis element.

The common traces identify \(\rho_T\otimes_T A\) with
\(\rho_\Delta\) integrally. One can check this directly.
The two idempotents of \(h\), with unit denominator
\(\omega-\omega^{-1}\), give free rank-one eigenspaces.
Complex conjugation \(c\) has square one, and both
\(\operatorname{tr}c\) and \(\operatorname{tr}(hc)\) vanish.
In that basis its diagonal entries vanish and its off-diagonal
entries are inverse units. Rescale one vector so \(c\) is the swap
matrix. The matrices \(1,h,c,hc\) form a basis of \(M_2(A)\) with
perfect trace pairing; only powers of \(2(\omega-\omega^{-1})\)
occur as denominators. Pairing \(\rho(g)\) with these matrices
recovers every entry from the shared traces.

The dual of (RCQ.6) is an isometric injection
\[
 A^\dagger=\operatorname{Hom}_{\mathcal O}(A,\mathcal O)
                   \hookrightarrow T^\dagger .
 \tag{RCQ.8}
\]
Every element of the unit ball of \(A\) has an integral lift in \(T\),
which proves the norm assertion. There is also an integral isometry
\[
 A\xrightarrow{\sim}A^\dagger,\qquad
 a\longmapsto\bigl(b\mapsto[1](ab)\bigr),
 \tag{RCQ.9}
\]
where \([1]\) takes the identity coefficient. The group-element
bases show that this Frobenius pairing has no \(|\Delta|^{-1}\).

\textit{2. Uniform bounds at \(p\).}
The restriction of \(\vartheta_\Delta\) to \(G_{K_p}\) is unramified.
If its Frobenius value is \(z\in\Delta\), then
\[
 \rho_{\Delta,p}\simeq\rho_{0,p}\otimes\nu,\qquad
             \nu(\operatorname{Frob}_p)=z^{1/2}.
 \tag{RCQ.10}
\]
The square root is unique in the finite odd-order group \(\Delta\),
and arithmetic \(\operatorname{Frob}_{K_p}\) is the square of
\(\operatorname{Frob}_p\).
Use AC-K33-ray-period, with the inverse orientation for a dual,
to identify this unramified twist inside the integral
\(\widetilde A^-\)-valued Kirillov ambient. The period and its
inverse are integral isometries.

Fix once the comparison for \(\rho_0\) between the Colmez--Wang
local factor and the Pan frame convention of AC-K33.
Colmez--Wang factors \(\rho_T\otimes_T\Pi(\rho_T^\bullet)\);
their action differs from the latter by
\(g\mapsto(g^{-1})^t\) (Remark 13.12).
This fixed comparison and its inverse have one bounded constant.
Unramified twisting commutes with the comparison; inverse transpose
inverts the determinant twist. Its values are units, so the
constant remains unchanged.

In the Pan convention the already checked simultaneous vectors are
\[
 \Phi_{j,d}(x)=1_{\mathbb Z_p^\times}(x)
       \sum_{k=0}^{j}(-1)^k\frac{j!}{(j-k)!}\,
        t^{-k}x^{-k}d\otimes e_1^k e_2^{j-k},
                     \quad d\in D_0.
 \tag{RCQ.11}
\]
Tensor this fixed construction with \(A^\dagger\), and make the
unit-period change (RCQ.10). In a group-element dual basis it is
a direct sum of identical fixed maps followed by an integral
isometry. Both upper and lower norm bounds are therefore independent
of \(\Delta\). The \(k=0\) coefficient retains the same injectivity
bound. For every \(g\in{\rm GL}_2(\mathbb Q_p)\), the unramified
factor is a power of the group-ring element \(z^{1/2}\); multiplication
by it permutes the group-element basis, so it and its inverse are
isometries. Together with the fixed unitary Banach action, this
proves the bound for all \(p\)-local \(G\)-translates, not only for
compact Kirillov restrictions.
An unramified determinant character is trivial on units and
multiplies \(U_p\) only by a unit, so the compact \(P_0\)-invariance,
inverse-diamond exponent and primitive-root calculation \(U_p=0\)
from AC-K33 remain valid. One depth \(a\) suffices for all \(\Delta\).

\textit{3. Integral away-prime vectors before splitting characters.}
At the finitely many ramified primes \(\ell\ne p\), apply
AC-K33-ray-Schwartz to the \(A\)-linear dual of
\(\rho_\Delta|_{G_{\mathbb Q_\ell}}\).
At unramified primes use the normalized spherical vector.
In the exact Colmez--Wang convention its values on
\(\operatorname{diag}(\ell^n,1)\), \(n\ge0\), are the complete
homogeneous polynomials in the Frobenius eigenvalues. They have
initial values \(1,\operatorname{tr}(r(\operatorname{Frob}_\ell))\)
and recurrence
\[
 a_n=\operatorname{tr}(r(\operatorname{Frob}_\ell))a_{n-1}
                  -\det(r(\operatorname{Frob}_\ell))a_{n-2};
\]
the negative values vanish. Thus the values lie in \(A\).
Iwasawa decomposition and the integral central character prove
integrality on every translate. This is the normalized formula
following Colmez--Wang Lemma 3.13.
Every factor has norm at most one and value one at the identity.
The same bounds hold for the product over arbitrarily many ramified
primes. These smooth vectors have a principal tame stabilizer
\(M_\Delta\).

\textit{4. The actual saturated completed-cohomology lattice.}
Put \(I=\ker(T\to A)\).
Rationally \(A[1/p]\) is a finite product of fields.
Every factor is finite Artin: it has no
\(\epsilon_p\)-ratio at \(\ell\ne p\), and at \(p\) is the fixed
irreducible induced representation times an unramified character.
All local minimum representations in Colmez--Wang \S3.1.2 are
therefore the full representations. The two inclusions of
Theorem 3.34(i) identify the rational \(I\)-annihilated global
representation with the full product of these factors.
This step is only rational; it asserts neither integral base change
nor connectedness.

Tensor the vectors of Steps 2 and 3 in that rational space and
view the coefficients through (RCQ.8).
Every global Whittaker value is integral after multiplication by
the one \(p^C\) from Step 2: Step 3 loses nothing, and the unitary
\(p\)-local action preserves the fixed bound.
It lies in the actual integral global representation by its defining
intersection
\[
 \begin{split}
 \Pi(\rho_T^\bullet)
 ={}&\iota^G\bigl(\Pi(\rho_T^\bullet)_{\rm ns}[1/p]\bigr)\\
 &{}\cap {\rm LC}\bigl(\widetilde G(\mathbb A_f^{(p)}),
              (\widetilde A^-\widehat\otimes\rho_T^*)^H\bigr).
 \end{split}
 \tag{RCQ.12}
\]
This is Colmez--Wang (3.30), using
\((\rho_T^\bullet)^\clubsuit=\rho_T^*\).
We have tested integrality in the intersection itself, rather than
bounded an index between a disconnected product of lattices and
its group-ring intersection.

The residual local representation is absolutely irreducible, so
the generic hypothesis of Colmez--Wang Theorem 13.11 holds.
That theorem puts these vectors in actual completed cohomology
integrally. Its full printed proof identifies the two images
isometrically, using Proposition 3.36 and Theorem 10.1.
The same fixed lower norm bound holds: evaluating all away factors
at the identity recovers the \(p\)-local vector, and (RCQ.8)--(RCQ.9)
are isometries. The scalar ambiguity in an identification of
\(\rho_T\otimes A\) with \(\rho_\Delta\) cancels against its dual,
exactly as in Colmez--Wang Remark 13.2(i).
We have obtained a uniformly bounded-below injection of the actual
finite coefficient module
\(\rho_\Delta\otimes_A A^\dagger\otimes_{\mathcal O}D_0^\circ\)
into the compact-support tower with \({\rm Sym}^j\)-coefficients.

\textit{5. The literal inverse tower.}
Apply the compact-support Cartan--Leray and duality construction of
AC-K33 with finite free \(\mathcal O\)-coefficients \(A\).
Tame descent divides by no group order: \(H_c^0\) of each open
modular curve vanishes, so the degree-one compact-support edge
identifies invariant classes. The \(p\)-group \(P_0\) and depth
\(a\) are fixed from Step 2. Deep \(p\)-level traces kill the
\(H_c^2\) colimit modulo every \(p^n\), just as in AC-K33.
Finite-coefficient duality and compact inverse limits identify the
bounded dual with the actual inverse trace tower in (RCQ.4).
All cyclotomic components of the full \(\mathbb Q\)-scheme are retained.

Restriction of bounded functionals to the injected finite-dimensional
space is onto by nonarchimedean Hahn--Banach. Its norm estimate uses
only the fixed lower bound from Step 4, not the rank of \(A\).
Restriction itself is \(T\)-linear and Galois-equivariant and hence
is \(A\)-linear after passing to that quotient; a chosen
Hahn--Banach extension need not be.
The \(\mathcal O\)-dual of
\(\rho_\Delta\otimes_A A^\dagger\) is canonically
\(\rho_\Delta^{\vee_A}\), integrally by (RCQ.9).
This proves (RCQ.4)--(RCQ.5).
The exact transpose calculation of AC-K33 supplies the inverse
diamond and \(U_p'\), unchanged by the unramified twist.
\end{proof}

\begin{remark}[Remaining class and norm questions]
The result supplies an actual group-ring quotient of cohomology
and a uniform bound for its coefficient projection.
The choices of unramified periods and tame compact vectors have
not been compared here with the geometric transition maps of the
full \(K\)-ray system. Hence no rank-two Euler norm identity over
\(K\) follows. Nor is the image of a preordinary Rankin class shown
to be nonzero: the determinant of its two local singular regulator
images remains the separate gate after AC-K33-image.
The integral argument uses no character-idempotent denominator and
does not invoke Colmez--Wang Conjecture 3.35.
\end{remark}
